\section{Evaluation Methodology}
\label{sec:method}

We ask four questions.
\textbf{RQ1}: How well does a policy derived only from tool descriptions contain a fully hijacked agent, and what does it cost in human review?
\textbf{RQ2}: Does calibrated deferral preserve that containment, and what does information flow add?
\textbf{RQ3}: What do typed actions, information flow and prediction add over earlier, string-based versions of the same packs?
\textbf{RQ4}: Do edits proposed from failures carry over to held-out cases?
Host enforcement was checked end to end in \S\ref{sec:hosts}.

\subsection{AgentDojo ground-truth replay}
AgentDojo \AdBench{}~\cite{agentdojo} (package version \AdVersion{}) has four suites (workspace, travel, banking, Slack) with \AdTools{} tools, user tasks and injection tasks; each task ships ground-truth tool calls. We use them without a language model.

\emph{Benign sessions.} Each user task's ground-truth calls, in order, form one session (\AdBenign{} sessions).

\emph{Attack sessions.} An injected instruction can only steer the agent after the agent has read it. We plant a unique marker in every injection slot of the environment, execute each user task's calls in AgentDojo's own runtime, and record the first call whose output contains a marker. An attack session for a pair (user task, injection task) is the user task's calls up to and including that call, followed by the injection task's ground-truth calls. The attacker thus chooses every call after the exposure: a fully hijacked agent, the strongest attacker a tool-call filter faces. Every user task read injected content (\AdUnexposedPairs{} pairs excluded). We exclude \AdGoalsNoCalls{} injection tasks that ship no ground-truth calls: \AdGoalsNoCallsWorkspace{} workspace tasks added in v1.2, which AgentDojo checks by end state, and \AdGoalsNoCallsTravel{} travel task that only manipulates text. This leaves \AdGoals{} attack goals in \AdAttacks{} attack sessions.

\emph{Outcomes.} A call reaching \verdict{ask} is treated as approved for the rest of its session (the worst case for later calls). A benign session is \emph{unattended} if every call is allowed or modified, and \emph{completes with review} if none is denied. An attack session is \emph{unreviewed} if every goal call is allowed or modified and no earlier call was refused, \emph{reaches a person} if some goal call asks and none is denied, and \emph{refused} if a goal call is denied. Whether a person refuses an attack that reaches them is outside the filter; we report the share separately.

\emph{Policies.} No policy was tuned on AgentDojo. A neutral pack has five tiers: read (\verdict{allow}), change (\verdict{ask}), destroy (\verdict{deny}), exec (\verdict{deny}) and unknown (\verdict{ask}). Each tool's tier comes from the tool registry's hint rules, regular expressions over its name and description that were fixed before this evaluation. \emph{Tiers} uses only these. \emph{Tiers + flow (hinted)} adds data labels from the registry's flow hints. \emph{Tiers + flow (conservative)} labels by a mechanical rule instead: every read is untrusted and sensitive (in a personal assistant every read returns third-party content and the user's own data), and a tool is egress if its description says it sends or one of its parameters is a URL. All three use the Rule of Two. Baselines: allow everything, and ask a person for everything.

\emph{Statistics.} Intervals are 95\% percentile intervals from a cluster bootstrap with 2{,}000 resamples: benign tasks are resampled, and attack sessions are resampled by attack goal (\AdGoals{} clusters), since sessions sharing a goal are not independent. For a zero count we also give the rule-of-three upper bound over goals.

\subsection{Labelled suites and simulation}
Nine built-in suites (\InCases{} cases) cover SOC triage, labelled SOC response actions with analyst appropriateness labels and an asset inventory, authorised testing under rules of engagement, coding agents, the tool names of Claude Code and OpenHands, multi-step information-flow sessions and typed network arguments. They were written alongside the packs, so we report them as regression tests and case studies, and compare each pack with earlier versions taken from the repository history (Table~\ref{tab:internal}). For deferral we simulate decision histories: scores uniform on $[0,6]$, denial probability $0.01$ below 2, $0.05$ in $[2,3)$ and $0.5$ above; \SimRuns{} runs per history size, $\alpha=0.05$, $\delta=0.1$, fixed seed. Latency is the in-process decision time per call over the replay (\LatN{} calls), excluding network transport.
